\section{Well-posed economic operators and implemented policies}
\label{supp:foundations119}
This section supplies sufficient conditions and proofs for the existence,
measurability, uniqueness, and stability obligations used in the paper.
These are standard dynamic-programming foundations, stated explicitly
rather than claimed as a new numerical contribution
\citep{ruszczynski2010risk,ruszczynski2014erratum}.  They apply directly
to the original investment primitives below.  A different recursive
utility, information filtration, or noncompact state model requires its
own corresponding hypotheses; it does not inherit these properties by
being described as a Bellman problem.

\subsection{A sufficient finite-horizon model}
Let $X$ and $\mathsf A$ be compact metric spaces.  At each finite date
$j<T$, let $A_j:X\rightrightarrows\mathsf A$ have nonempty compact
values and be both upper and lower hemicontinuous.  Its graph is denoted
$G_j$.  Let the innovation space $\mathsf Z_j$ be compact metric with
fixed Borel probability $P_j$.  The original transition
$F_j:G_j\times\mathsf Z_j\to X$ is jointly continuous and invariant
in $X$.  The cost $c_j:G_j\to\mathbb R$ and terminal payoff
$g:X\to\mathbb R$ are continuous.  The finite discounts
$\beta_j\ge0$ need not be smaller than one.

Let $\rho_j:L^\infty(P_j)\to\mathbb R$ be finite, monotone, cash
invariant, and normalized by $\rho_j(0)=0$.  Without normalization,
the identical original operator is obtained using
$\bar\rho_j=\rho_j-\rho_j(0)$ and
$\bar c_j=c_j+\beta_j\rho_j(0)$; thus the finite deterministic
constant remains in the original stage account.  For implemented Borel policies we also
require the following parameter-measurability property:
\begin{equation}\label{eq:kernelmeas119}
 (s,z)\mapsto w(s,z)\text{ is bounded and Borel}\quad\Longrightarrow
 \quad s\mapsto\rho_j(w(s,\cdot))\text{ is Borel}
\end{equation}
for every standard Borel parameter space and bounded Borel $w$.
This property is verified below for the original expectation and AVaR
primitives; it is not deduced from cash invariance alone.

The filtration is generated by the successive original innovations,
with their stipulated fixed conditional laws, and the initial state is
observed.  An admissible history policy is a nonanticipative Borel action
rule on this history, feasible at its current original state.  Its
criterion is the \emph{specified conditional nested recursion}: at a
history, apply $\rho_j$ to the next-innovation section of its
continuation.  This specifies the time-consistent original problem.
Static AVaR of total cost and a randomization device with a different
observation convention are different criteria and are not inserted into
this policy class without being declared.

\begin{proposition}[Finite-horizon existence, selection, and validation]
\label{prop:finitefoundations119}
Under these conditions, the original Bellman recursion has one bounded
continuous value sequence $V_T=g$,
\[
 V_j(x)=\min_{a\in A_j(x)}
 \{c_j(x,a)+\beta_j\rho_j(V_{j+1}(F_j(x,a,\cdot)))\}.
\]
Each minimum is attained and admits a Borel selector.  The resulting
Markov policy is optimal among all admissible history policies for the
stipulated recursion.  Every implemented Borel feasible policy has a
well-defined bounded Borel value sequence; its value is at least $V_j$
at every date and original history with state $x$.

Uniqueness here concerns the value sequence.  An optimizer need not be
unique.  If each action objective has a unique minimizer, that minimizer
is continuous in $x$.  In particular, a strongly convex objective on
a compact convex feasible action set has a unique minimizer.

For the same primitives, two bounded continuous terminal continuations satisfy
\begin{equation}\label{eq:terminalstability119}
 \|V_j-\widetilde V_j\|_\infty
 \le\left(\prod_{k=j}^{T-1}\beta_k\right)
                  \|g-\widetilde g\|_\infty.
\end{equation}
For a fixed implemented policy the same bound holds for bounded Borel
terminal continuations, whose recursions are well-defined by
\eqref{eq:kernelmeas119}.  Additive
stage-cost discrepancies of at most $e_k$ add their dated discounted
sum to the right side.  Changes in a law, transition, constraints,
or preference require their own operator discrepancy; they are not
covered by this same-primitives statement.
\end{proposition}

\begin{proof}
If $\|v-w\|_\infty\le h$, then $w-h\le v\le w+h$.
Monotonicity and cash invariance give
\[
 |\rho_j(v)-\rho_j(w)|\le h,
 \qquad |\rho_j(v)|\le\|v\|_\infty.
\]
This also applies to equivalence classes under $P_j$; a stronger
pointwise bound is sufficient.  The feasible graph is closed in a
compact product and hence compact.  For continuous $v$, joint
continuity on $G_j\times\mathsf Z_j$ implies uniform continuity.
Consequently $(x,a)\mapsto v(F_j(x,a,\cdot))$ is continuous in
the supremum norm.  The displayed nonexpansiveness proves continuity
of its risk evaluation, and thus of the full action objective.

Compactness of $A_j(x)$ gives an attained minimum.  To prove value
continuity without hiding a selection assumption, let $x_n\to x$.
Along a subsequence attaining the lower limit of the values, compactness
gives a convergent subsequence of minimizers $a_n\to a$.
Upper hemicontinuity gives $a\in A_j(x)$, and continuity of the
objective implies that the lower limit is at least $V_j(x)$.
Conversely, choose a minimizer $a$ at $x$.  Lower hemicontinuity
provides $a_n\in A_j(x_n)$ with $a_n\to a$; evaluating these
feasible actions shows that the upper limit is at most $V_j(x)$.
Thus the value is continuous.  Backward induction from $g$ proves
existence and continuity of the complete value sequence; the finite
recursion also proves its uniqueness.  The norm estimate
$\|V_j\|_\infty\le\|c_j\|_\infty+
\beta_j\|V_{j+1}\|_\infty$ proves boundedness.

Here is an explicit measurable-selector argument.  The argmin graph
is closed, because both the objective and value are continuous, and
its fibers are nonempty compact.  For each $n$, choose a finite cover
of $\mathsf A$ by closed balls of radius $2^{-n}$ with fixed
centers.  In a fixed ordering choose the first ball meeting the argmin
fiber, and at each successive level choose the first ball meeting its
intersection with the previously chosen balls.  All these intersections
remain nonempty compact and are nested.  The event that the argmin
fiber meets a fixed finite intersection of closed balls is closed:
it is the projection onto compact $X$ of a closed subset of
$X\times\mathsf A$.  Induction over the finite-valued previous
choices therefore makes each new chosen center a Borel function of
$x$.  The nested sets have diameters at most $2^{1-n}$, so their
intersection contains exactly one action.  The chosen centers converge
to that action; their pointwise limit is Borel and belongs to the
argmin fiber.  This constructs the required selector.  When the
argmin is already unique, the same subsequence argument used for
value continuity shows continuity of the selector.

For an arbitrary implemented history policy, composition of Borel
histories, actions, and the original transition gives a bounded Borel
next-history kernel.  Property~\eqref{eq:kernelmeas119} then
constructs its bounded Borel value backward from $g$.  Suppose at
the next date that the policy continuation is at least
$V_{j+1}$ at its next original state.  Applying the same conditional
risk at the current history preserves this inequality.  The resulting
current cost is at least the action objective at its feasible chosen
action, hence at least $V_j$.  Backward induction proves policy
domination, without restricting the comparison to Markov policies.
Using the Borel minimizing selector at every date makes every
inequality an equality, which proves its optimality.

Finally, the preceding nonexpansiveness yields
$\|\mathcal T_jv-\mathcal T_jw\|_\infty
\le\beta_j\|v-w\|_\infty$, by comparing infima over the same
actions.  A fixed policy has the identical inequality without the
infimum.  Backward induction proves
\eqref{eq:terminalstability119}.  When the costs differ by at
most $e_j$, the one-step estimate has the extra term $e_j$;
iterating it gives the dated sum.
\end{proof}

\paragraph{The exact continuity and measurability obligations.}
Compact innovations can be replaced by essential-uniform continuity
of $v(F_j(x,a,\cdot))$ on the declared continuation class.  Pointwise
continuity alone is insufficient for an arbitrary monetary risk.  For
example, let $Z$ have an unbounded exponential law,
$F(x,z)=1-e^{-xz}$ for $x\in[0,1]$, terminal value $g(y)=y$,
and $\rho=\operatorname{ess\,sup}$.  All these primitives are
continuous and bounded where required, but
$\rho(g(F(0,Z)))=0$ and $\rho(g(F(x,Z)))=1$ for every $x>0$.
Expectation has stronger dominated-continuity properties, but those
cannot be assigned to every monotone cash-invariant risk.

For the original fixed-law AVaR and bounded jointly Borel $w$, the
formula
\[
 \rho_\alpha(w(s,\cdot))=
 \inf_{\eta\in\mathbb Q}
 \left\{\eta+\frac1{1-\alpha}
                  \int(w(s,z)-\eta)_+\,P(dz)\right\}
\]
proves \eqref{eq:kernelmeas119}.  Each integral is Borel: the
claim holds on measurable rectangles, extends to indicators of
product-Borel sets by a monotone-class argument, and then to
nonnegative Borel functions by monotone simple approximation.
The infimum is countable and equals the real-threshold infimum
because the displayed expression is continuous in $\eta$.
Expectation is the same measurability argument without an infimum.
Thus the needed property is established for both original preference
cases, while remaining explicit for a different risk primitive.

\subsection{The original investment domain and unique constrained actions}
\begin{proposition}[Original primitive closure and actor regularity]
\label{prop:investmentfoundations119}
For the original even-dimensional investment primitives, every
feasible state-action-innovation triple satisfies
\begin{equation}\label{eq:invariance119}
 F(x,a,z)\in[1/32,27/32]^d\subset[0,1]^d.
\end{equation}
The capacity correspondence $[0,A(x)]$ is continuous, compact,
and nonempty, with $1/8\le A(x)\le1/4$.  Costs are continuous
and nonnegative, with valid uniform bounds
$S\le49/32$, $g\le37/16$, and $S+h\le103/64$.
Hence the finite-horizon existence and policy theorem applies for
the original expectation and nested AVaR criteria.

For the two-date AVaR model with $\beta=15/16$, $V_1$ is strongly convex in state
with modulus $(4-5\beta/4)/d$, and the original root and last-date
action objectives are strongly convex in action with modulus two.
They have unique continuous constrained optimal actors.  The
last-date $V_1$ and the root action objective are continuously
differentiable on the relative state interior, with the corresponding
one-sided boundary interpretation.

For either date where $Q_j$ is continuously differentiable, write
$a_j^*(x)$ for its constrained minimizer and
\[
 \lambda_j(x)=\boldsymbol1_{\{a_j^*=A(x)\}}
                  [-\partial_aQ_j(x,a_j^*)]_+.
\]
Then the value gradient is
\begin{equation}\label{eq:constrainedenvelope119}
 \nabla V_j(x)=\nabla_xQ_j(x,a_j^*(x))
                         -\lambda_j(x)\nabla A(x).
\end{equation}
This formula is not a derivative of a rounded deployed actor.
Interior actor derivatives require the additional local second
derivatives and strict-interiority conditions stated in the retained
compiler theorem.
\end{proposition}

\begin{proof}
For the state part of the $j$th transition coordinate put
$p=x_j$, $q=x_{j+1}$.  The function
$p/2+q/8+p(1-q)/16$ is nondecreasing in both arguments on
$[0,1]^2$, because its two partial derivatives are
$1/2+(1-q)/16$ and $1/8-p/16$.  Its range lies in $[0,5/8]$.
Since $0\le b_ja\le1/8$ and $|z|\le1/32$, adding the
constant $1/16$ proves \eqref{eq:invariance119}.  The affine
capacity gives the stated correspondence properties.  The inequalities
$|x_j-5/8|\le5/8$, $|x_j-x_{j+1}|\le1$, and
$(1/2-2\bar x)_+\le1/2$ give the stated state-cost bounds.
Moreover $h(a)\le h(1/4)=5/64$, which gives the full stage
bound.  Polynomial transitions, squared positive parts, and the fixed
compact innovation interval have the required continuity.  The
previous measurability proof applies to the risk primitive.

For completeness, the inherited joint curvature bound follows directly
from the original primitives.  On each smooth region,
$\nabla^2S\succeq(4/d)I$, $\nabla^2g\succeq0$, and
$|\partial_jg|\le10/d$.  The only nonzero second derivative
of the transition is the mixed derivative of
$-x_jx_{j+1}/16$.  In a joint state-action direction $(v,w)$,
the Hessian chain rule therefore gives
\[
 D^2(g\circ F)[(v,w)]^2
 \ge-\frac18\sum_j\partial_jg(F)v_jv_{j+1}
 \ge-\frac5{4d}\|v\|_2^2.
\]
The omitted Jacobian--Hessian--Jacobian term is nonnegative in
the full joint direction, including its state-action mixed terms.
Together with $h''\ge2$, this proves for every fixed shock
the lower curvature
\[
 \frac{4-5\beta/4}{d}\|v\|_2^2+2w^2.
\]
The squared-hinge derivatives match at the junction, so integrating
the almost-everywhere second-derivative inequality along a segment
extends it across the junction.  Integration under any feasible
AVaR density preserves the bound.  Subtracting its common quadratic
term leaves convex functions; their supremum is convex, so the risk
supremum preserves the same bound.  It follows that the joint
objective has the state term
$(4-5\beta/4)\|x-y\|_2^2/(2d)$ in its strong-convexity
inequality.  Interpolating two optimal state-action pairs preserves
feasibility because the capacity is affine.  Taking the value at that
feasible interpolation and dropping the nonnegative action-distance
term proves strong convexity of $V_1$ with the announced modulus.
This supplies the state convexity used in the active-tail theorem,
rather than assuming it for a fitted continuation.

For each fixed root state, the successor is affine in its scalar
action.  Convexity of $V_1$, followed by monotonicity and convexity
of AVaR, makes $a\mapsto\rho_\alpha(V_1(F(x,a,\cdot)))$
convex.  Adding $h''\ge2$ proves root-action strong convexity;
the same property of the last-date objective follows from the
original primitive curvature.  The preceding compact theorem now
gives existence, uniqueness, and continuity of both optimal actors.

The retained exact last-date formula is $C^1$ in state and action:
the positive-part square and its explicit tail polynomial have matching
values and first derivatives at their junctions.  Normalize the action
by $a=A(x)v$, with $v\in[0,1]$.  The normalized objective has
a unique continuous minimizer.  Comparing the value increment at the
old optimizer and at the new optimizer, and using uniform first-order
expansion on this fixed compact action interval, gives
\[
 \nabla V_1=\nabla_xQ_1+v^*\partial_aQ_1\nabla A.
\]
At a lower endpoint $v^*=0$, at an interior optimum
$\partial_aQ_1=0$, and at capacity
$v^*=1$, $\partial_aQ_1\le0$.  This proves
\eqref{eq:constrainedenvelope119} and continuity of the value
gradient.  In particular the multiplier tends to zero at a switch
between an interior and a binding-capacity solution.

For the root objective, invariance puts every successor strictly
inside the state cube, where $\nabla V_1$ is continuous and
bounded on the compact successor range.  Along the innovation,
$V_1(F(x,a,z))$ is strongly convex with modulus $181/64$.
Every one of its level sets has at most two points, so its
uniform-law distribution has no atoms.  Continuity of the function
on the interval makes its distribution function strictly increasing
between its minimum and maximum: the preimage of every nonempty
open value interval in that range contains a nonempty open innovation
interval.  Its quantile is therefore unique at the declared risk
levels.  The differentiability conditions of the retained AVaR
envelope argument apply, with bounded action/state derivatives.
Equivalently, the active central interval has a unique continuous
minimizer and the value envelope differentiates the integral with
that interval held fixed.  Thus, writing $E_{x,a}$ for the optimal
two-tail set,
\[
 \partial_aQ_0(x,a)=h'(a)+
 \frac{\beta}{2R(1-\alpha)}
       \int_{E_{x,a}} b^{\mathsf T}\nabla V_1(F(x,a,z))\,dz.
\]
The analogous state derivative has $D_xF^{\mathsf T}\nabla V_1$
and $\nabla S$.  Continuity of the integrands and tail endpoints
gives continuity of these derivatives.  Applying the same normalized
action envelope proves \eqref{eq:constrainedenvelope119} for
$V_0$.  Endpoint actions are included; no global differentiability
of the actor across active-set changes is asserted.
\end{proof}

\paragraph{Observation, feasibility, and terminal-policy transfer.}
The implemented terminal controller observes a downward $40$-bit state
$y\le x$ with $\|x-y\|_1\le d2^{-40}$ and returns an action
feasible at $y$.  Since $A(y)\le A(x)$, it remains feasible at
the original state.  Its certified action gap at $y$ is at most
$2^{-21}$.  The primitive fixed-action objective and the optimal
value both have state Lipschitz constant at most $15/d$, so its
original-state gap is at most
\[
 2^{-21}+30\,2^{-40}<2^{-20}.
\]
The fixed-action modulus follows from
$\|\nabla S\|_\infty\le8/d$, the transition column-sum bound
$11/16$, $\|\nabla g\|_\infty\le10/d$, and
$8+\beta(11/16)10<15$.  The value modulus is the retained
capacity-aware bound.  The controller's rounding and exact fallback
must actually enforce its premise; an optimizer return code alone
does not do so.  This explains why its allowance remains valid when
the new method changes the locations of original-law integration
queries while retaining the same deployed controller.

\subsection{Residual validation and the optional infinite-horizon case}
For the finite-horizon one-sided sandwich, suppose
$f_j-\mathcal T_jf_{j+1}\le u_j$,
$\mathcal T_j^\pi f_{j+1}-f_j\le v_j$, and
$\ell_T\le f_T-g\le u_T$.  Monotonicity and cash invariance
give separately
\[
 f_0-V_0\le\sum_{j<T}B_ju_j+B_Tu_T,
 \qquad
 J_0^\pi-f_0\le\sum_{j<T}B_jv_j-B_T\ell_T.
\]
Each follows by backward induction, starting from its corresponding
terminal inequality.  Adding them proves the retained complete-policy
bound, and feasibility supplies $J_0^\pi\ge V_0$.  This validates
an NBO's own attained continuation without assuming optimal-continuation
training labels.  Date-specific constant shifts of $f_j$ cancel in
the paired residuals and in $u_T-\ell_T$.

For a stationary infinite-horizon model with the same compact and
kernel conditions, bounded cost, and $0\le\beta<1$, the Bellman
map is a $\beta$-contraction on the complete space $C(X)$.
The finite-horizon continuity proof shows that it maps this space
into itself.  Thus it has a unique continuous fixed point $V$.
For any Borel feasible stationary policy, the policy map is a
$\beta$-contraction on bounded Borel functions by
\eqref{eq:kernelmeas119}; it has a unique bounded Borel fixed
point $J^\pi$.  A Borel minimizer of the continuous action
objective at $V$ satisfies $T_{\pi^*}V=TV=V$, so uniqueness
of its policy fixed point gives $J^{\pi^*}=V$.
For another policy, $V\le T_\pi V$; monotone iteration and
contraction give $V\le J^\pi$.

The infinite-history criterion is well-defined as the limit of the
same finite nested recursions: the omitted cost after $n$ dates is
bounded uniformly by $\|c\|_\infty\beta^n/(1-\beta)$.
Passing the finite-horizon comparison to this limit proves optimality
against the same admissible history-policy class.  This argument does
not assert that infimizing arbitrary Borel action objectives preserves
Borel measurability.  The more general policy-iteration theorem in
the retained supplement explicitly assumes that additional preservation
and selection property when it applies a Bellman infimum to possibly
discontinuous policy values.

\subsection{The critic and actor computational graph}
For a smooth parameterized critic and differentiable feasible actor,
the derivative must correspond to the objective actually optimized.
For example, at a fixed original state write
\[
 r_\theta=c(x,a_\theta)+
     \beta\rho(f_{+,\theta}(F(x,a_\theta,Z)))-f_\theta(x).
\]
Suppose differentiation is justified by integrable local domination,
and either $\rho$ is expectation or it is AVaR at an atomless unique
quantile, as in the retained tail-gradient theorem.  Let $q_\theta$
be its maximizing tail density, with $q_\theta=1$ for expectation.
The chain and risk-envelope rules give
\begin{align*}
 \nabla_\theta r_\theta={}&
 (D_\theta a_\theta)^{\mathsf T}
 \left[\nabla_ac+
  \beta\mathbb E_{q_\theta}
          (D_aF)^{\mathsf T}\nabla_xf_{+,\theta}(F)\right]\\
 &+\beta\mathbb E_{q_\theta}\nabla_\theta f_{+,\theta}(F)
       -\nabla_\theta f_\theta(x).
\end{align*}
The last expectation is a partial parameter derivative at a fixed
successor.  The actor contribution is written separately to avoid
double counting.  A risk-envelope derivative holds the maximizing
density fixed in this formula; its validity comes from the stated
regularity assumptions, not from differentiating a sorting permutation.
The half-squared-residual objective
$\mathcal L(\theta)=\tfrac12\mathbb E[r_\theta^2]$ has gradient
$\mathbb E[r_\theta\nabla_\theta r_\theta]$ when its own
domination condition holds.  Detaching the future target defines a
different frozen-target gradient or a semigradient algorithm.  It
does not establish that this full original coupled gradient was used,
and neither formula proves optimizer convergence.

At a projection or rounding discontinuity these differentiability
premises may fail.  Piecewise training derivatives can be specified
as an algorithmic choice, but the global economic certificate still
uses the actual feasible action and its original-law inequalities.
This is consistent with the own-future constructive theorem, whose
validity does not depend on a successful differentiable optimizer.

\section{Stopping boundaries and the continuous-time recursion}
\label{supp:bsdeboundaries119}
The retained Brownian Markov bridge is a separate classical recursive
utility verification theorem
\citep{PardouxPeng1990,ElKarouiPengQuenez1997}.  Its forward
coefficients and driver are understood to be jointly Borel in their
displayed deterministic arguments, so composition with each admissible
progressively measurable control yields progressively measurable
coefficients.  The forward state-Lipschitz constants and driver
$(y,z)$-Lipschitz constants are uniform in the control.  Terminal
payoff and the driver at $(y,z)=(0,0)$ are bounded.  A feedback
must define its own well-posed closed-loop state process.  These
conditions make the retained existence proof precise; they do not
establish existence for a fully coupled FBSDE or an arbitrary
discontinuous feedback.

\begin{proposition}[Stopped original-BSDE verification]
\label{prop:stoppedbsde119}
Let $\tau^a\le T$ be a specified stopping time for each admissible
control, such as exit from a declared domain capped at $T$, with
bounded Borel original stopping payoff $G(\tau^a,X_{\tau^a})$.
Before stopping use exactly the driver and forward coefficients of
the retained BSDE bridge.  Let bounded $u\in C^{1,2}$ have the bounded
derivatives required for stopped It\^o calculus and continuous
values on the declared stopping boundary.  Suppose
\[
 \mathcal H^a u\ge-e_-(s)\ \text{for every admissible action},
 \qquad\mathcal H^\pi u\le e_+(s)
\]
on the unstopped domain, and at every possible stopping location
and at the terminal part of the boundary suppose
\[
 -\eta_-\le G(s,x)-u(s,x)\le\eta_+.
\]
Here $e_\pm$ are nonnegative deterministic integrable functions
and $\eta_\pm\ge0$.  The stopped BSDE has a unique solution
in $S^2\times H^2$, and
\[
 0\le Y_0^\pi-\inf_aY_0^a
 \le e^{L_yT}(\eta_-+\eta_+)
       +\int_0^T e^{L_ys}(e_-(s)+e_+(s))\,ds.
\]
The stopping rule and boundary payoff are parts of the original
economic criterion.  A smaller interior residual does not remove
their errors or permit them to be changed.
\end{proposition}

\begin{proof}
Embed the stopped equation on $[0,T]$ with terminal variable
$\xi^a=G(\tau^a,X_{\tau^a})$ and driver
$\boldsymbol1_{\{s<\tau^a\}}f(s,X_s,y,z,a_s)$.
It is progressively measurable, has the same uniform Lipschitz
constants, and has square-integrable terminal data and bounded
driver at zero.  The weighted Picard argument in the retained proof
therefore gives a unique $S^2\times H^2$ solution.  After
$\tau^a$ it equals $\xi^a$ with $Z=0$; before stopping it
is exactly the declared stopped equation.

Apply stopped It\^o calculus to $u(s\wedge\tau^a,
X_{s\wedge\tau^a})$.  Subtracting gives the same linearized
error BSDE as in the retained proof, with its coefficients set to
zero after $\tau^a$, terminal error
$\xi^a-u(\tau^a,X_{\tau^a})$, and residual driver
$\boldsymbol1_{\{s<\tau^a\}}\mathcal H^{a_s}u$.
Its measurable coefficients satisfy $|A_s|\le L_y$ and
$\|\vartheta_s\|\le L_z$.  The bounded Girsanov coefficient
has an exponential density with all finite moments.  The same
Cauchy--Schwarz and $H^1$ martingale argument used in the
retained verification proof therefore justifies expectation of the
stopped integration-factor identity,
\[
 Y_0^a-u(0,x)=\mathbb E^{\vartheta}\left[
 e^{\int_0^{\tau^a}A_vdv}
       (G-u)(\tau^a,X_{\tau^a})
 +\int_0^{\tau^a}e^{\int_0^sA_vdv}
                \mathcal H^{a_s}u(s,X_s)\,ds\right].
\]
The factor is positive and at most $e^{L_ys}$.
Applying the lower residual and boundary bounds gives a lower
bound for every admissible control.  Applying the upper bounds
gives an upper bound for the implemented policy.  Since
$\tau^a\le T$, the deterministic integrals and terminal factors
are bounded by those displayed in the proposition.  Taking the
infimum of the lower bounds and subtracting proves the result.
\end{proof}

\paragraph{Why the boundary and graph are material.}
For a deterministic state starting at $1/2$ in $(0,1)$, use either
of the two constant controls $a=\pm1$ in $dX_s=a\,ds$,
stop on exit before $T=1$,
and assign boundary costs zero at the left endpoint and one at the
right.  With zero driver and $u=0$, every interior residual is
zero, but the right-exit policy has regret one.  The stopping-boundary
error exactly supplies the missing term.  This is a verification
counterexample, not a second numerical application.

The BSDE graph is $z_u=\sigma^{\mathsf T}\nabla u$ for the
same control and the original driver.  Its $C^{1,2}$ requirement
is substantive.  A piecewise-affine ReLU critic is generally not
$C^2$ and cannot be inserted into this classical formula by declaring
its second derivative zero away from its kinks; a different weak,
viscosity, or smoothing argument with the required error account
would first be needed.  The finite-horizon ReLU Bellman construction
uses order and Lipschitz arguments and does not rely on this
continuous-time smooth verification theorem.

Finally, a general recursive utility is not automatically cash
invariant.  Even the globally Lipschitz BSDE driver $f(s,y,z)=\lambda y$
gives $Y_0=e^{\lambda T}\mathbb E\xi$, so adding a constant
$c$ to the terminal criterion changes $Y_0$ by
$e^{\lambda T}c$, not generally by $c$.  The continuous-time
bridge correctly uses its $e^{L_yT}$ stability factors.  It does
not inherit the monetary discrete-risk contraction or identify an
arbitrary NDU/FBSDE with the original investment AVaR objective.
